% !TeX root = ../../../Main.tex
\chapter{Literature Review}

\section{How to use the nomenclature}

The nomenclature entries must all be listed in the \texttt{nomenclature.tex} file located in \texttt{src/parts/frontmatter/}. This file organizes symbols into four categories: Acronyms (A), Latin letters (L), Greek letters (G), and Subscripts/Indices (S).

\subsection*{Entry Format}
Each nomenclature entry follows this syntax:
\begin{verbatim}
\nomenclature[Category]{Symbol}{Description}
\end{verbatim}

Where \texttt{Category} is one of: \texttt{A} (Acronyms), \texttt{L} (Latin letters), \texttt{G} (Greek letters), or \texttt{S} (Subscripts/Indices).

\subsection*{Examples}

\begin{enumerate}
    \item \textbf{Acronyms (Category A)}:
    \begin{verbatim}
\nomenclature[A]{KOS}{Koordinatensystem}
    \end{verbatim}
    This creates an entry for "KOS" with the description "Koordinatensystem".
    
    \item \textbf{Latin letters (Category L)}:
    \begin{verbatim}
\nomenclature[L]{$\boldsymbol{r}$}{\nomenunit{-}Vektor vom Lasthaken zum Schwerpunkt}
    \end{verbatim}
    The \texttt{\textbackslash nomenunit\{\}} command specifies the unit.
    
    \item \textbf{Greek letters (Category G)}:
    \begin{verbatim}
\nomenclature[G]{$\Omega$}{\nomenunit{\unitfrac{rad}{s}} Drehgeschwindigkeit}
    \end{verbatim}
    This adds the Greek letter Omega with its unit and description.
    
    \item \textbf{Subscripts/Indices (Category S)}:
    \begin{verbatim}
\nomenclature[S]{$_a(\cdot)$}{Aufhängepunktfestes Koordinatensystem}
    \end{verbatim}
    This defines the subscript 'a' used in equations throughout the document.
\end{enumerate}

\subsection*{Important Notes}
\begin{itemize}
    \item Always use \texttt{\textbackslash nomenunit\{\}} to specify units, with a hyphen (-) for dimensionless quantities.
    \item The entries will automatically appear in the appropriate subsection (Abbreviations, Latin letters, etc.) based on their category.
    \item After adding or modifying entries, you may need to run \texttt{makeindex} to update the nomenclature listing.
\end{itemize}

% ----------------------------------------------------------------------------------------------
\section{How to cite from the bibliography file}

Your references are stored in a bibliography file (with extension \texttt{.bib}). In this document template, the bibliography file is located at \texttt{src/bib/references.bib}. Open this file to see how bibliography entries are structured.

A typical bibliography entry looks like this:

\begin{verbatim}    
@BOOK{Johnson1994,
  title = {Helicopter theory},
  publisher = {Dover Publications},
  year = {1994},
  author = {Johnson, Wayne},
  address = {Mineola, NY}
}
\end{verbatim}

The entry starts with the document type (\texttt{@article}, \texttt{@book}, \texttt{@inproceedings}, etc.) followed by a citation key (\texttt{Johnson1994}) that you'll use to cite this source in your document.


There are two main approaches to building your bibliography:

\subsubsection*{Manual Entry}
For small bibliographies, you can create entries by hand. This gives you complete control but requires attention to detail. Follow the formats shown in your \texttt{references.bib} file.

\subsubsection*{Automatic Entry (Recommended)}
For larger bibliographies, use reference management tools to automatically add entries:

\begin{itemize}
    \item \textbf{Reference managers}: Tools like \href{https://www.zotero.org/}{Zotero}, \href{https://www.mendeley.com/}{Mendeley}, or \href{https://endnote.com/}{EndNote} can export collections as BibTeX files.
    \item \textbf{DOI to BibTeX}: Use websites like \url{https://doi2bib.org} to convert DOIs to BibTeX entries.
    \item \textbf{CrossRef}: Visit \url{https://www.crossref.org} to look up and export citations.
\end{itemize}


Always check automatically generated entries for consistency and completeness. Pay special attention to:
\begin{itemize}
    \item Proper capitalization (protect with braces, e.g., \texttt{{NASA}})
    \item Special characters and accents (use proper \LaTeX\ encoding)
    \item Complete page numbers and DOIs when available
\end{itemize}

Once you have entries in your bibliography file, you can cite them using various commands. This template uses \texttt{biblatex} with the \texttt{numeric} style, which provides flexible citation options:


\begin{table}[H]
    \centering
    \caption{Common \texttt{biblatex} citation commands}
    \label{tbl:citation-commands}
    \begin{tabular}{lll}
        \toprule
        \textbf{Command} & \textbf{Example Usage} & \textbf{Result} \\
        \midrule
        \texttt{\textbackslash cite\{key\}} & \texttt{\textbackslash cite\{Fichter2009\}} & \cite{Fichter2009} \\
        \texttt{\textbackslash textcite\{key\}} & \texttt{\textbackslash textcite\{Fichter2009\}} & Textual: \textcite{Fichter2009} \\
        \texttt{\textbackslash citeauthor\{key\}} & \texttt{\textbackslash citeauthor\{Fichter2009\}} & Only authors: \citeauthor{Fichter2009} \\
        \texttt{\textbackslash citeyear\{key\}} & \texttt{\textbackslash citeyear\{Fichter2009\}} & Only year: \citeyear{Fichter2009} \\
        \texttt{\textbackslash cite[42]\{key\}} & \texttt{\textbackslash cite[42]\{Fichter2009\}} & With page number: \cite[42]{Fichter2009} \\
        \texttt{\textbackslash cite[see][]\{key\}} & \texttt{\textbackslash cite[see][42]\{Fichter2009\}} & With pre/post notes: \cite[see][42]{Fichter2009} \\
        \texttt{\textbackslash cite\{key1,key2\}} & \texttt{\textbackslash cite\{Fichter2009,Johnson1994\}} & Multiple: \cite{Fichter2009,Johnson1994} \\
        \bottomrule
    \end{tabular}
\end{table}
